\subsection{Locally convex final topologies}

PROPOSITION 5. --- Let E be a vector space, and \((\mathrm{F}_{\alpha})_{\alpha\in\mathbf{A}}\) be a family of topological vector spaces and for each \(\alpha\in A\) , let \(g_{\alpha}\) be a linear mapping of \(F_{\alpha}\) in E.

\begin{enumerate}
\def\labelenumi{(\roman{enumi})}
\item
  Denote by \(\mathfrak{B}\) the family of absorbent, symmetric convex subsets \(\mathrm{V}\) of \(\mathrm{E}\) such that \(g_{\alpha}^{-1}(\mathrm{V})\) is a neighbourhood of 0 in \(\mathrm{F}_{\alpha}\) for every \(\alpha\) ; the family \(\mathfrak{B}\) is a fundamental system of neighbourhoods of 0 in \(\mathrm{E}\) for a topology \(\mathcal{T}\) that is compatible with the vector space structure.
\item
  A linear mapping \(f\) of \(\mathbf{E}\) in a locally convex space \(\mathbf{G}\) (resp. a semi-norm \(p\) on \(\mathbf{E}\) ) is continuous for \(\mathcal{T}\) if and only if, for every index \(\alpha, f \circ g_{\alpha}\) (resp. \(p \circ g_{\alpha}\) ) is continuous in \(\mathrm{F}_{\alpha}\) .
\item
  The topology \(\mathcal{T}\) is the finest of the locally convex topologies on \(\mathbf{E}\) for which the \(g_{\alpha}\) are continuous.
\end{enumerate}

Further, the topology \(\mathcal{T}\) is the only locally convex topology on \(\mathbf{E}\) that satisfies condition (ii) for linear mappings (resp. for the semi-norms).

As \(\mathfrak{B}\) is a filter base invariant under homotheties of ratio \(>0\) , the assertion (i) follows immediately from II, p.~23, prop. 1. By the definition of \(\mathfrak{B}\) , the topology \(\mathcal{T}\) is the finest of locally convex topologies on E making the \(g_{\alpha}\) continuous; whence (iii). Finally, it is clear that if \(f\) is continuous, so is \(f \circ g_{\alpha}\) ; conversely if the \(f \circ g_{\alpha}\) are continuous for every \(\alpha\) , then for each symmetric convex neighbourhood W of 0 in G, the set \(g_{\alpha}^{-1}(f^{-1}(\mathbf{W}))\) is a neighbourhood of 0 in \(\mathbf{F}_{\alpha}\) for each \(\alpha\) . Now \(f^{-1}(\mathbf{W})\) is absorbent, symmetric and convex thus \(f^{-1}(\mathbf{W})\) is a neighbourhood of 0 in \(\mathcal{T}\) , and \(f\) is continuous. Similarly if \(p\) is a semi-norm on E such that \(p \circ g_{\alpha}\) is continuous for every \(\alpha\) , and if U is the set of points \(x \in \mathbf{E}\) such that \(p(x) < 1\) , then, for every \(\alpha\) , the set \(g_{\alpha}^{-1}(\mathbf{U})\) is a convex neighbourhood of 0 in \(\mathbf{E}_{\alpha}\) that is symmetric and absorbent; thus U is a neighbourhood of 0 in E and \(p\) is continuous (II, p.~2, prop. 1).

The last statement follows from S, IV, § 2.5, criterion CST 18.

We say that \(\mathcal{T}\) is the locally convex final topology of the family of topologies \(\mathcal{T}_{\alpha}\) of the \(F_{\alpha}\) , for the family of linear mappings \(g_{\alpha}\) .

It may happen that \(\mathcal{T}\) is not the finest of the topologies on E compatible with its vector space structure and making the \(f_{\alpha}\) continuous (II, p.~75, exerc. 15; see also II, p.~75, exerc. 14).

In the most important case \(E = \sum_{\alpha \in A} g_{\alpha}(F_{\alpha})\) , we get a fundamental system of neighbourhoods of 0 for T as follows; for each \(\alpha \in A\) , let \(V_{\alpha}\) be a symmetric neighbourhood of 0 for \(T_{\alpha}\) , form the union of the \(g_{\alpha}(V_{\alpha})\) for \(\alpha \in A\) and denote the convex envelope in E of this union by \(\Gamma((g_{\alpha}(V_{\alpha})))\) ; since every element of E is of the form \(\sum_{\alpha \in J} x_{\alpha}\) , where J is a finite subset of I and \(x_{\alpha} \in g_{\alpha}(F_{\alpha})\) , it is immediate that \(\Gamma((g_{\alpha}(V_{\alpha})))\) is an absorbent symmetric convex set in E (each of the \(V_{\alpha}\) is absorbent in \(F_{\alpha}\) ); as \(\Gamma((g_{\alpha}(V_{\alpha})))\) contains all the \(g_{\alpha}(V_{\alpha})\) , it is a neighbourhood of 0 for T. On the other hand, it is clear that for every symmetric convex neighbourhood V of 0 for T, we have \(V \supset \Gamma((V \cap g_{\alpha}(F_{\alpha})))\) , from which our assertion follows.

COROLLARY 1. --- With the notations of prop. 5, let H be a set of linear mappings of E in the locally convex space G. Suppose that E is the sum of its subspaces \(g_{\alpha}(F_{\alpha})\) ; then H is equicontinuous for T, if, and only if, for every \(\alpha\) , the set \(f \circ g_{\alpha}\) where f varies in H, is equicontinuous in \(F_{\alpha}\) .

Remembering I, p.~9, prop. 6 the argument is similar to that of (ii) prop. 5. Let W be a symmetric convex neighbourhood of 0 in G and note that if the set \(f \circ g_{\alpha}\) , where \(f \in \mathbf{H}\) is equicontinuous, then the intersection \(\bigcap_{f \in \mathbf{H}} g_{\alpha}^{-1}(f^{-1}(\mathbf{W}))\) is a symmetric convex neighbourhood of 0 in \(\mathbf{F}_{\alpha}\) . As this intersection is the same as \(g_{\alpha}^{-1}\big(\bigcap_{f \in \mathbf{H}} f^{-1}(\mathbf{W})\big)\) and the set \(\bigcap_{f \in \mathbf{H}} f^{-1}(\mathbf{W})\) is symmetric and convex, everything depends on showing that it is also absorbent. Now, by hypothesis, every \(x \in \mathbf{E}\) can be written as \(\sum_{i=1}^{n} g_{\alpha_i}(z_{\alpha_i})\) , where \(z_{\alpha_i} \in \mathrm{F}_{\alpha_i}\) . To show that there exists \(\lambda > 0\) such that \(f(\lambda x) \in \mathbf{W}\) for all \(f \in \mathbf{H}\) , it is sufficient to consider the case \(x = g_{\alpha}(z_{\alpha})\) with \(z_{\alpha} \in \mathrm{F}_{\alpha}\) (since we can pass to the general case by replacing \(\mathbf{W}\) by \(\mathbf{W}/n\) ). But this case follows from the fact that \(g_{\alpha}^{-1}\big(\bigcap_{f \in \mathbf{H}} f^{-1}(\mathbf{W})\big)\) is a neighbourhood of 0 in \(\mathrm{F}_{\alpha}\) .

COROLLARY 2. --- Let \((\mathbf{J}_{\lambda})_{\lambda \in \mathbf{L}}\) be a partition of the index set A. Let \((\mathbf{G}_{\alpha})_{\alpha \in \mathbf{A}}\) be a family of locally convex spaces and \((\mathbf{F}_{\lambda})_{\lambda \in \mathbf{L}}\) be a family of vector spaces. For each \(\lambda \in \mathbf{L}\) , let \(h_{\lambda}\) be a linear mapping of \(\mathbf{F}_{\lambda}\) in a vector space E; for each \(\lambda \in \mathbf{L}\) and \(\alpha \in \mathbf{J}_{\lambda}\) , let \(g_{\lambda \alpha}\) be a linear mapping of \(\mathbf{G}_{\alpha}\) in \(\mathbf{F}_{\lambda}\) . Write \(f_{\alpha} = h_{\lambda} \circ g_{\lambda \alpha}\) . Suppose that each \(\mathbf{F}_{\lambda}\) carries the finest locally convex topology that makes the \(g_{\lambda \alpha} (\alpha \in \mathbf{J}_{\lambda})\) continuous. Then, the finest locally convex topology on E that makes the \(f_{\alpha}\) continuous, is identical with the finest locally convex topology making the \(h_{\lambda}\) continuous.

This is a particular case of S, IV, § 2.5 criterion CST 19, and can also be proved directly using prop. 5.

Examples of locally convex final topologies.

\paragraph{I. Quotient space.}

Let M be a subspace of the locally convex space F, and \(\phi\) be the canonical mapping of F on F/M. As the quotient topology on F/M is locally convex and is the finest of all the topologies (locally convex or not) which make \(\phi\) continuous, it is also the locally convex final topology for the family consisting of the single mapping \(\phi\) .

\paragraph{II. Inductive limits of locally convex spaces.}

Let A be an ordered set directed to the right and let \((\mathrm{E}_{\alpha}, f_{\beta\alpha})\) be an inductive system of vector spaces relative to the set A (A, II, § 6.2); let \(E = \varinjlim E_{\alpha}\) and let \(f_{\alpha}: E_{\alpha} \to E\) be the canonical linear mapping for each \(\alpha \in A\) . Suppose that each \(E_{\alpha}\) carries a locally convex topology \(T_{\alpha}\) , and further suppose that for \(\alpha \leqslant \beta\) , the mapping \(f_{\beta\alpha}: E_{\alpha} \to E_{\beta}\) is continuous. Then we say that the locally convex final topology T of the family \((\mathcal{T}_{\alpha})\) relative to the linear mappings \(f_{\alpha}\) (resp. the space E carrying the topology T) is the inductive limit of the family \((\mathcal{T}_{\alpha})\) (resp. the inductive limit space of the system \((\mathrm{E}_{\alpha}, f_{\beta\alpha})\) , or simply of the locally convex spaces \(E_{\alpha}\) ). Recall that E is the union of the vector subspaces \(f_{\alpha}(\mathrm{E}_{\alpha})\) and that when \(\alpha \leqslant \beta\) , we have \(f_{\alpha}(\mathrm{E}_{\alpha}) \subset f_{\beta}(\mathrm{E}_{\beta})\) ; if we endow \(f_{\alpha}(\mathrm{E}_{\alpha})\) with the final topology for the mapping \(f_{\alpha}\) (which is the same as identifying \(f_{\alpha}(\mathrm{E}_{\alpha})\) with the quotient space \(E_{\alpha}/f_{\alpha}^{-1}(0)\) ), the topology T is also the final topology of the family of the topologies of the \(f_{\alpha}(\mathrm{E}_{\alpha})\) , relative to the canonical injections (II, cor. 2 above). Further, the continuity of \(f_{\beta\alpha}\) for \(\alpha \leqslant \beta\) implies that the canonical injection \(j_{\beta\alpha}: f_{\alpha}(\mathrm{E}_{\alpha}) \to f_{\beta}(\mathrm{E}_{\beta})\) is continuous, so that E is also the inductive limit of \(f_{\alpha}(\mathrm{E}_{\alpha})\) carrying the preceding topologies relative to the injection \(j_{\beta\alpha}\) .

Example. --- Let X be a locally compact space and \(E = \mathcal{K}(X; \mathbf{R})\) the vector space of finite continuous real valued functions defined over X with compact support. For every compact subset K of X, let \(E_{K}\) be the vector subspace of E formed by those functions \(f \in E\) which are such that \(x \notin K \Rightarrow f(x) = 0\) . Denote by \(T_{K}\) the topology induced on \(E_{K}\) and by \(T_{u}\) the topology of uniform convergence on X. The inductive limit T of the topologies \(T_{K}\) is finer than \(T_{u}\) ; we can show that if X is paracompact and not compact, then T is strictly finer than \(T_{u}\) (cf.~INT, III, 2nd ed., § 1.8). The importance of T lies in the fact that the linear forms on E that are continuous in T are precisely the real measures on X (INT, III, 2nd., § 1.3).

Remark. --- In the last example, the topology induced by \(\mathcal{T}\) on \(\mathrm{E}_{\mathbf{K}}\) is identical with \(\mathcal{T}_{\mathbf{K}}\) , since by definition it is coarser than \(\mathcal{T}_{\mathbf{K}}\) and, since \(\mathcal{T}\) is finer than \(\mathcal{T}_{u}\) , the topology induced by \(\mathcal{T}\) on \(\mathrm{E}_{\mathbf{K}}\) is finer than that induced by \(\mathcal{T}_{u}\) , that is to say \(\mathcal{T}_{\mathbf{K}}\) .

This reasoning generalises immediately to an inductive limit of locally convex topologies \((\mathcal{T}_{\alpha})\) when there is a locally convex topology \(\mathcal{T}'\) on E such that \(\mathcal{T}_{\alpha}\) is the topology induced on \(\mathrm{E}_{\alpha}\) by \(\mathcal{T}'\) .

More generally one can ask, when we suppose that \(E_{\beta} \subset E_{\alpha}\) and \(T_{\beta}\) is the topology induced by \(T_{\alpha}\) , under what circumstances T induces \(T_{\alpha}\) on each of the \(E_{\alpha}\) . In general this is not so (II, p.~80, exerc. 26); but we shall see in the Nos following two important situations where this does occur.

